\chapter{Empirical Study}
\label{chap:empirical}

This chapter presents the findings of the empirical part of the study.
It includes the survey data gathered from forty medium-sized and large
Latvian enterprises with formal PMO structures, along with findings
from one structured expert interview. The chapter is organized into
five sections: sample characteristics, descriptive statistics,
hypothesis testing, expert interview findings, and a chapter summary.

\section{Sample Characteristics}
\label{sec:sample}

A total of forty valid survey responses were collected, satisfying the
30--50 response target specified in the research methodology. The
respondents represented organizations from a variety of industries and
sizes, providing a diverse and representative sample for the analysis.

\subsection{Organization Size}
Over half of the respondents were from medium-sized enterprises.
Specifically, 52.5\% (n=21) reported having 50--249 employees, 35\%
(n=14) had 250--999 employees, and 12.5\% (n=5) represented large
organizations with 1000 or more employees. This distribution reflects
the actual business landscape in Latvia, which is characterized by
a dominance of medium-sized companies.

\subsection{Industry Distribution}
The sample covered seven different industry sectors. The largest group
was Information Technology at 37.5\% (n=15), which is not surprising
given that IT companies tend to have more formal PMO structures. Finance
and Banking, Retail and Commerce, and Telecommunications each accounted
for 12.5\% (n=5). Public Sector and Other industries each represented
10\% (n=4), while Manufacturing made up 5\% (n=2).

\subsection{PMO Maturity and Scope}
Respondents represented organizations with varying levels of PMO
establishment and project portfolio sizes. This diversity strengthens
the reliability of the findings by ensuring that the results are not
skewed toward any single type of organization.

\section{Descriptive Statistics}
\label{sec:descriptive}

\subsection{Integration Maturity Level Distribution}
One of the central questions of this study was to determine at which
integration maturity level most Latvian PMO environments currently
operate. The results are presented in Table~\ref{tab:imm-distribution}.

{\singlespacing
\begin{table}[h]
  \caption{Distribution of IS Integration Maturity Levels}
  \label{tab:imm-distribution}
  \centering
  \begin{tabular}{p{0.08\linewidth}p{0.20\linewidth}p{0.15\linewidth}p{0.15\linewidth}}
    \toprule
    \textbf{Level} & \textbf{Name} & \textbf{Frequency} & \textbf{Percentage}\\
    \midrule
    1 & Initial & 6 & 15.0\%\\
    2 & Developing & 13 & 32.5\%\\
    3 & Defined & 15 & 37.5\%\\
    4 & Optimized & 6 & 15.0\%\\
    \midrule
    & \textbf{Total} & \textbf{40} & \textbf{100\%}\\
    \bottomrule
  \end{tabular}
\end{table}
}

The results clearly show that the majority of PMO environments (70\%,
n=28) operate at intermediate maturity levels — either Level 2
(Developing) or Level 3 (Defined). Only 15\% (n=6) have reached the
highest level of integration maturity, and another 15\% (n=6) remain
at the most basic level. This pattern suggests that while many Latvian
enterprises have moved beyond entirely manual integration, full
platform-level integration remains uncommon.

\subsection{Descriptive Statistics for Key Constructs}
Table~\ref{tab:descriptive} shows the mean scores and standard
deviations for the three dependent variables — reporting automation,
data quality, and governance standardization — measured using
five-point Likert scales.

{\singlespacing
\begin{table}[h]
  \caption{Descriptive Statistics for Key Constructs}
  \label{tab:descriptive}
  \centering
  \begin{tabular}{p{0.35\linewidth}p{0.10\linewidth}p{0.10\linewidth}p{0.10\linewidth}p{0.10\linewidth}p{0.10\linewidth}}
    \toprule
    \textbf{Construct} & \textbf{N} & \textbf{Mean} & \textbf{Std Dev}
    & \textbf{Min} & \textbf{Max}\\
    \midrule
    Reporting Automation & 40 & 3.38 & 0.85 & 1.00 & 5.00\\
    Data Quality & 40 & 3.32 & 0.73 & 2.00 & 5.00\\
    Governance Standardization & 40 & 3.64 & 0.71 & 1.00 & 5.00\\
    \bottomrule
  \end{tabular}
\end{table}
}

All three constructs show moderate mean scores, suggesting that the
sampled organizations have achieved a moderate but not yet high level
of performance in these areas. Governance standardization recorded
the highest mean (M=3.64), while data quality had the lowest (M=3.32).
The standard deviations indicate a reasonable spread of responses,
reflecting genuine variation in IS integration practices across the
sample.

\section{Hypothesis Testing}
\label{sec:hypotheses}

Before conducting correlation analysis, the Shapiro-Wilk normality
test was applied to all variables. The results indicated that the
data did not follow a normal distribution, which led to the use of
Spearman rank correlation — a non-parametric method suitable for
ordinal and non-normally distributed data \cite{field2018}.
Statistical significance was assessed at the p $<$ 0.05 level.

\subsection{Hypothesis 1}
\textbf{H1: The majority of PMOs in medium and large Latvian
enterprises operate at intermediate integration maturity levels
(Levels 2 or 3) rather than at the most advanced level.}

As shown in Table~\ref{tab:imm-distribution}, 70\% of respondents
(n=28) reported operating at Level 2 or Level 3, while only 15\%
(n=6) had reached Level 4. This clearly confirms that most Latvian
PMO environments are at intermediate rather than advanced integration
maturity levels. \textbf{Hypothesis 1 is supported.}

\subsection{Hypothesis 2}
\textbf{H2: PMOs that use API or platform-based integration report
higher levels of reporting automation compared to those that rely
on manual data handling.}

Spearman correlation analysis showed a statistically significant
positive relationship between integration maturity level and
reporting automation (r=0.405, p=0.010). This means that
organizations with higher IS integration maturity consistently
reported higher levels of automated reporting.
\textbf{Hypothesis 2 is supported.}

\subsection{Hypothesis 3}
\textbf{H3: Higher integration maturity is associated with better
data accuracy and faster access to project information.}

A statistically significant positive correlation was found between
integration maturity level and data quality (r=0.321, p=0.043).
Although this is the weakest of the three correlations, it remains
statistically significant and confirms that better integrated systems
are associated with higher perceived data accuracy and timeliness.
\textbf{Hypothesis 3 is supported.}

\subsection{Hypothesis 4}
\textbf{H4: PMOs that use centralized BI or dashboard tools tend
to have more standardized governance practices.}

The correlation analysis revealed a statistically significant
positive relationship between integration maturity level and
governance standardization (r=0.403, p=0.010). Organizations
with higher integration maturity demonstrated notably stronger
governance standardization practices.
\textbf{Hypothesis 4 is supported.}

{\singlespacing
\begin{table}[h]
  \caption{Summary of Hypothesis Testing Results}
  \label{tab:hypothesis-results}
  \centering
  \begin{tabular}{p{0.05\linewidth}p{0.38\linewidth}p{0.10\linewidth}
  p{0.10\linewidth}p{0.15\linewidth}}
    \toprule
    \textbf{H} & \textbf{Hypothesis} & \textbf{r} & \textbf{p}
    & \textbf{Result}\\
    \midrule
    H1 & Most PMOs at Levels 2--3 & -- & -- & Supported\\
    H2 & Maturity $\to$ Reporting Automation & 0.405 & 0.010 & Supported\\
    H3 & Maturity $\to$ Data Quality & 0.321 & 0.043 & Supported\\
    H4 & Maturity $\to$ Governance & 0.403 & 0.010 & Supported\\
    \bottomrule
  \end{tabular}
\end{table}
}

\section{Expert Interview Findings}
\label{sec:interviews}

One structured expert interview was conducted to add qualitative
depth to the survey findings. The participant is a professional
with experience in PMO environments. In line with the ethical
principles outlined in Chapter~2, the participant is referred
to anonymously as Expert 1.

\subsection{Expert 1 — Project Management Professional}

Expert 1 is a skilled project management professional working in a
PMO environment. Their organization operates a hybrid-integrated IS
ecosystem, utilizing tools such as ServiceNow or Planview for
strategic portfolio management, Jira for agile execution, and
Microsoft Project Online for waterfall-style schedules. These systems
are connected through APIs and middleware solutions, placing the
organization at a high integration maturity level (Level 3--4 of
the IMM).

Despite this high level of integration, Expert 1 identified three
recurring challenges. First, data latency — although systems are
connected, short delays between execution tools and the portfolio
layer can cause stakeholders to view slightly outdated reports.
Second, semantic misalignment — different departments define concepts
such as progress in different ways, requiring ongoing reconciliation
between technical and financial systems. Third, legacy resistance —
older ERP systems used by the Finance department are difficult to
connect with modern cloud-based tools, sometimes requiring manual
data exports as a workaround.

Regarding reporting, Expert 1 described a 70/30 split in favour of
automation. Standard KPIs, budget consumption, and resource
utilization are pulled automatically into Power BI and Tableau
dashboards, while the remaining 30\% — primarily executive summaries
and qualitative narratives — still requires human input.

On data quality, Expert 1 rated satisfaction as 4 out of 5.
Timeliness had improved significantly due to automated syncing,
but accuracy remains dependent on how consistently project managers
update their data. To address this, the organization introduced a
Data Health Score system to monitor data hygiene across projects.

Expert 1 confirmed that a centralized Power BI Command Center had
been a major step forward for governance standardization. It enforced
universal KPI definitions across all projects, eliminated subjective
reporting, and made project performance visible to all executives
at the same time.

For recommendations, Expert 1 suggested that Latvian enterprises
should consider modular, low-code platforms such as Monday.com or
ClickUp, which offer strong integration capabilities at a lower cost
than enterprise-level tools. Expert 1 also highlighted the need for
standardized connectors between local Latvian accounting software
and global PM tools, and stressed the importance of data literacy
training for Latvian project leaders.

\section{Chapter Summary}
\label{sec:chap3summary}

This chapter presented the empirical findings of the study. Based
on 40 valid survey responses from medium and large Latvian
enterprises, the results showed that 70\% of PMO environments
operate at intermediate integration maturity levels (Levels 2 and 3),
which confirms Hypothesis 1. Statistically significant positive
correlations were found between integration maturity and all three
dependent variables: reporting automation (r=0.405, p=0.010),
data quality (r=0.321, p=0.043), and governance standardization
(r=0.403, p=0.010), confirming Hypotheses 2, 3, and 4. The expert
interview with Expert 1 provided rich qualitative support for
these findings, particularly highlighting the value of centralized
BI systems for governance standardization and the ongoing challenge
of legacy system integration. The overall findings are discussed
and interpreted in Chapter~4.
